\documentclass[10pt,a4paper]{amsart}
\usepackage[margin=19mm]{geometry}
\usepackage{amsmath,amssymb,mathtools}
\usepackage[scr=rsfs,cal=euler]{mathalfa}
\usepackage{xfrac}
\usepackage[unicode,hidelinks,bookmarks=false]{hyperref}
\newcommand{\ie}{i.e.\ }
\newcommand{\cf}{cf.\ }
\newcommand{\resp}{respectively\ }
\newcommand{\Subsection}{\subsection}
\providecommand{\nobreakdash}{\nobreak\hbox{-}}
\setlength{\emergencystretch}{2em}
\multlinegap0pt
\usepackage{xcolor}
\newtheorem{defi}{Definition}
\newtheorem{prop}{Proposition}
\newcommand{\CC}{\mathcal{C}}
\newcommand{\Ho}{\mathrm{H}}
\newcommand{\Uc}{\mathcal{U}}
\newcommand{\co}{\colon}
\title{Baues-Wirsching cohomology and\\ \v{S}varc genus in small categories}
\author{Isaac Carcac\'ia-Campos, Enrique Mac\'ias-Virg\'os, David Mosquera-Lois}
\date{Source: arXiv:2508.00727v2 (CC BY 4.0)}
\begin{document}
\maketitle

\noindent\emph{Source and licence.} Baues-Wirsching cohomology and\\ \v{S}varc genus in small categories. Version arXiv:2508.00727v2 (CC BY 4.0), distributed by arXiv under the Creative Commons Attribution 4.0 International licence, http://creativecommons.org/licenses/by/4.0/, as recorded in the arXiv licence metadata for this exact version and shown on the abstract page https://arxiv.org/abs/2508.00727v2. Full licence notice: \texttt{licenses/arxiv-2508.00727-CC-BY-4.0.txt}.\par\smallskip

\noindent\textit{Editorial scope.} Counted unit: the article's prop rel-cup (source0.tex lines 530--532). The article's complete original proof of it is delivered with it (source0.tex lines 534--560, one \texttt{proof} environment of 27 lines). No mathematics is written, repaired or re-derived: the statement and the proof are the article's own lines, in the article's own notation, and no retained line is substituted or re-flowed.  Retained uncounted as the source-local context the proof needs: lines 473--475; lines 481--483; lines 518--528.  Nothing was omitted from any retained span, so the package carries no omission record.  No retained line contains a citation, so the wrapper carries no bibliography and no bibliography entry.  No retained line contains a figure environment, an image-inclusion command or drawing code, so no image is delivered and the package declares no asset dependency.  Editorial formatting only, in addition: an AMS \texttt{amsart} preamble that restates the article's own declarations and macros used by the retained lines, namely the environments \texttt{defi}, \texttt{prop} on separate counters, together with the standard packages those lines need.  The retained environments print in this document as Definition 1, Proposition 1 in order of appearance, whereas the article prints the same items with the numbering of its own full document.  The complete native source of the article as distributed in the arXiv e-print for this exact version is bundled unchanged as \texttt{sources/182567/source0.tex}.
    \begin{equation}\label{eq:cup_product}
        \delta(f \smile g)= (\delta f) \smile g + (-1)^n f \smile (\delta g)
    \end{equation}

\begin{equation}\label{Product_Induction}
    \delta(f_1 \smile \dots \smile f_n)= \sum_{i=1}^{n} (-1)^{N_i} f_1 \smile \dots \smile (\delta f_i) \smile \dots \smile f_n
\end{equation}

\begin{defi}\label{defi:cup_product_relative}
    Let $\CC$ be a category, $\Uc$ be a subcategory and $\{\Uc_0, \dots, \Uc_n\}$ a geometric cover of $\Uc=\Uc_0 \cup \ldots \cup \Uc_n$. Let $D$ be a natural system in $\CC$ with an endopairing $\mu \colon (D,D) \rightarrow D$. We define \begin{equation}\label{eq:relative_cup_product_chain_complexes}
    \smile\co F^{p_0}(\CC,\Uc_0;D) \otimes \cdots \otimes F^{p_n}(\CC,\Uc_n;D) \rightarrow F^{p}(\CC,\Uc;D)
    \end{equation}
    (where $p=p_0+\dots+p_n$) by $$(f_0 \smile \dots \smile f_n)(\lambda_1,\dots,\lambda_{p})=f_0(\lambda_1,\dots,\lambda_{p_{0}}) \cdots f_n(\lambda_{P_{n}},\dots,\lambda_{p})$$
    where $P_{i}$ is the sum $p_0+\cdots + p_{i-1}$.
    This induces a \emph{relative cup product}
\begin{equation}\label{eq:relative_cup_product_cohomology}
        \smile\co \Ho^{p_0}(\CC,\Uc_0;D) \otimes \cdots \otimes \Ho^{p_n}(\CC,\Uc_n;D) \rightarrow \Ho^{p}(\CC,\Uc;D).
    \end{equation}  
\end{defi}

\begin{prop}\label{rel-cup}\label{VARIOS}
    The products defined in Equations (\ref{eq:relative_cup_product_chain_complexes}) and (\ref{eq:relative_cup_product_cohomology}) in Definition \ref{defi:cup_product_relative} are well-defined. 
\end{prop}

\begin{proof}
    The morphism (\ref{eq:relative_cup_product_chain_complexes})
    is well defined since we have a geometric cover. Indeed, the cochain $f_0 \smile \dots \smile f_n$ vanishes in every chain in $\Uc$. By definition, for any $p$-chain $(\lambda_1,\ldots,\lambda_{p})$ we have that
    $$(f_0 \smile \dots \smile f_n)(\lambda_1,\dots,\lambda_{p})=f_0(\lambda_1,\dots,\lambda_{p_0})  \cdots f_n(\lambda_{P_n+1},\dots,\lambda_{p}).$$ But since the chain $(\lambda_1,\ldots,\lambda_{p})$ is in $\Uc$ and we have a geometric cover, the chain belongs to some $\Uc_i$. Hence the subchain $(\lambda_{P_i+1},\dots,\lambda_{P_{i+1}})$ is in $\Uc_i$, so $f_i(\lambda_{P_i+1},\dots,\lambda_{P_{i+1}})=0$ and the cup-product also vanishes.
    
    Now we will see that morphism \eqref{eq:relative_cup_product_cohomology} is  well-defined.  Formula \eqref{Product_Induction} ensures after a tedious computation that if we have two sequences $f_0, \dots, f_n$ and $f'_0, \dots, f'_n$ such that $\{f_i\}=\{f'_i\}$ then $$\{f_0 \smile \dots \smile f_n\}=\{f'_0 \smile \dots \smile f'_n\}.$$ 
    In the case of two elements the proof goes as follows. 
    
    If we have four cochains $f, f' \in F^n(\CC,\Uc_0;D)$ and $g, g' \in F^n(\CC,\Uc_1;D)$ such that $\{f\}=\{f'\}$ and $\{g\}=\{g'\}$, we will prove that $\{f \smile g\}=\{f'\smile g'\}$. 
    
    As we know $f=f'+\delta h_1$ and $g=g'+\delta h_2$ with $h_1 \in F^{n-1}(\CC,\Uc_0;D)$ and $h_2 \in F^{m-1}(\CC,\Uc_1;D)$, by using the fact that $\delta$ is a derivation (Formula \eqref{eq:cup_product}) we have 
    \begin{align*}
         &(f \smile g)(\lambda_1,\ldots,\lambda_n, \lambda_{n+1}, \ldots \lambda_{n+m})\\
         &=f(\lambda_1,\ldots,\lambda_n) \cdot g(\lambda_{n+1},\ldots,\lambda_{n+m}) \\
         &=(f'+\delta h_1)(\lambda_1,\ldots,\lambda_n) \cdot (g'+\delta h_2)(\lambda_{n+1},\ldots,\lambda_{n+m}) \\
         &= f'(\lambda_1,\ldots,\lambda_n)\cdot g' (\lambda_{n+1},\ldots,\lambda_{n+m})\\ &\quad + f'(\lambda_1,\ldots,\lambda_n) \cdot \delta h_2 (\lambda_{n+1},\ldots,\lambda_{n+m}) \\ &\quad + \delta h_1(\lambda_1,\ldots,\lambda_n) \cdot g' (\lambda_{n+1},\ldots,\lambda_{n+m})\\
        & \quad+ \delta h_1 (\lambda_1,\ldots,\lambda_n) \cdot \delta h_2(\lambda_{n+1},\ldots,\lambda_{n+m}) \\
          &= (f' \smile g' + f' \smile \delta h_2 + \delta h_1 \smile g' + \delta h_1 \smile h_2) (\lambda_1,\ldots\lambda_{n+m}).
    \end{align*}
    But we also have that 
    \begin{enumerate}
        \item $\delta(f \smile h_2)=\delta f \smile h_2+(-1)^{n}f \smile \delta h_2=(-1)^{n}f \smile \delta h_2$, since $\delta f=0$;
        \item  $\delta(h_1 \smile g)=\delta h_1 \smile g+(-1)^{n-1} h_1 \smile \delta g= \delta h_1 \smile g$, since $\delta g=0$;
        \item $\delta(h_1 \smile \delta h_2)= \delta h_1 \smile \delta h_2 +(-1)^{n-1} h_1 \smile \delta  \delta h_2 =\delta h_1 \smile \delta h_2$, since $\delta \delta =0$.
    \end{enumerate}
    Therefore $$f \smile g= f'\smile g' +\delta((-1)^{n}f \smile h_2 + h_1 \smile g + h_1 \smile \delta h_2).\qedhere$$
\end{proof}

\end{document}
